\section{Mechanisation and Extraction}\label{sec:mech}
The scheme of Section~\ref{sec:schulze} is formalised in four files on
top of the library: \rocqfile[Frontend/Schulze.v]{src/Frontend/Schulze.v} (the ballot, its voter and
verifier, aggregation), \rocqfile[Backend/SchulzeTally.v]{src/Backend/SchulzeTally.v} (the Schulze
method over $\mathbb{Z}$, the certificate, its prover and verifier, the
count), \rocqfile[Utility/SchnorrOrder.v]{src/Utility/SchnorrOrder.v} (the order of elements of the
Schnorr group) and \rocqfile[Examples/SchulzeIns.v]{src/Examples/SchulzeIns.v} (the instantiation and
the running example). All of it is generic in the vector space of
Section~\ref{schn_proof}, and nothing is admitted. The main theorems
are:
\begin{itemize}\setlength{\itemsep}{0pt}
  \item \emph{Honest ballots are accepted.} A ballot made by an honest
  voter, with positions in $[0,m]$, passes the verifier, both in the
  interactive protocol and with Fiat--Shamir (\rocq{ballot_valid},
  \rocq{nizk_ballot_valid}).
  \item \emph{Accepted ballots are well formed.} From two accepting runs
  that share the announcements but differ in the challenge, one can
  extract a position vector $s$ such that every bit ciphertext encrypts
  $[s_i < s_j]$. So an accepted ballot encrypts one ranking
  (\rocq{ballot_sound}, Listing~\ref{lst:ballot-sound}).
  \item \emph{Ballots reveal nothing.} Every ballot transcript has exactly
  the distribution of a simulator that does not know the ranking
  (\rocq{ballot_zk}).
  \item \emph{Aggregation is correct.} The product of the accepted
  ballots' ciphertexts encrypts the pairwise counts $D$
  (\rocq{aggregate_margins_enc}, \rocq{cdec_aggregate}).
  \item \emph{The certificate pins down the winner.} The conditions of
  part (iii) hold exactly for the min-max closure, the closure is unique,
  and a candidate that passes part (iv) is a Schulze winner
  (\rocq{cert_ok_is_closure}, \rocq{closure_unique}, \rocq{cert_winner}).
  \item \emph{Honest trustees succeed.} Trustees who run the method
  correctly produce a certificate that the verifier accepts, both
  interactively and with Fiat--Shamir (\rocq{verify_tally_complete},
  \rocq{nizk_tally_complete}).
  \item \emph{Accepted certificates are correct.} By rewinding, an
  accepted certificate names a Schulze winner of the encrypted counts
  (\rocq{verify_tally_sound}).
  \item \emph{Certificates reveal only the paths.} Every certificate
  transcript has exactly the distribution of a simulator that gets the
  fresh ciphertexts, the challenges, and the witness paths
  (\rocq{tally_zk}).
  \item \emph{End to end.} Accepted ballots and an accepted certificate
  on their aggregate name a Schulze winner of the rankings the voters
  cast, and so does a successful run of the count
  (\rocq{schulze_end_to_end}, \rocq{count_sound}).
\end{itemize}

The batched ballot of Section~\ref{sec:batched} is formalised in
\rocqfile[Frontend/SchulzeSimple.v]{src/Frontend/SchulzeSimple.v}, on top
of the library's Bayer--Groth zero argument and a new generic
$\Sigma$-protocol for the preimage of a linear map
(\rocqfile[Crypto/LinearSigma.v]{src/Crypto/LinearSigma.v}). The new
protocol comes with the usual three theorems: completeness, exact special
soundness, and perfect zero-knowledge in the permutation form
(\rocq{lin_completeness}, \rocq{lin_soundness}, \rocq{lin_shvzk}). For
the batched ballot itself:
\begin{itemize}\setlength{\itemsep}{0pt}
  \item \emph{The derived key is correct.} Committing to $s$ under
  $\mathbf{g}'$ gives the same group element as committing to $Ms$ under
  $\mathbf{g}$ (\rocq{com_gd}).
  \item \emph{Honest ballots are accepted}, interactively and with
  Fiat--Shamir (\rocq{ballot3_valid}, \rocq{nizk_ballot3_valid}).
  \item \emph{Soundness and zero-knowledge} are the theorems
  \rocq{ballot3_sound} and \rocq{ballot3_zk} of Section~\ref{sec:batched}.
  \item \emph{End to end.} Either the commitment key is broken, or the
  certified candidate is a Schulze winner of the committed positions
  (\rocq{schulze_end_to_end3}, in
  \rocqfile[Backend/SchulzeSimpleTally.v]{src/Backend/SchulzeSimpleTally.v}).
\end{itemize}

Listing~\ref{lst:ballot-sound} shows \rocq{ballot_sound}. It
concludes with a position vector $s$, $s_i \le m$, such that every pair
ciphertext has plaintext $g^{[s_i < s_j]}$: whatever the voter's device
did, the bits of an accepted ballot are the bits of one ranking. The hypotheses are the key relation $h = g^x$, the
two transcripts, and two side conditions that the library's generic
statements cannot discharge: \rocq{exp_injective}~\texttt{g}, that $g$ has order
$q$ so that $g^a = g^b$ implies $a = b$ in $\mathbb{Z}_q$, and
\rocq{natF_injective}~\texttt{(4m+2)}, that the field has characteristic above
$4m+2$ so that the small integers of the membership proofs are distinct
field elements. \rocq{verify_tally_sound} has the same shape for the
certificate, with the bound $2N^2 + 3N$. It is proved by running the
library's extractors (\rocq{generalised_or_sigma_soundness_main} and
the decryption-proof extractor) on every proof of the certificate and
decoding the extracted plaintexts through the decryption homomorphism.
\rocq{schulze_end_to_end} composes the two with
\rocq{cdec_aggregate}, which says that the plaintext of the aggregate
is the sum of the ballots' bits, and with the lemma of
Section~\ref{sec:tally}.

\begin{lstlisting}[frame=single, language=Coq, caption={Ballot soundness ({\rocq[\texttt{Frontend/Schulze.v}]{ballot_sound}})},
  label={lst:ballot-sound},captionpos=t, basicstyle=\ttfamily\scriptsize,
  abovecaptionskip=-\medskipamount]
Theorem ballot_sound : forall (n : nat) (x : F) (g h : G)
  (b b' : @ballot (S n)), h = g ^ x -> 
  exp_injective g -> 
  natF_injective (4 * S n + 2) -> 
  same_announcements b b' -> 
  different_challenges b b' ->
  verify_ballot g h b = true -> 
  verify_ballot g h b' = true ->
  exists s : Vector.t nat (S n),
    (forall i, Vector.nth s i <= S n) /\
    (forall i j, plaintext_bit x b i j =
      g ^ (if Nat.ltb (Vector.nth s i) (Vector.nth s j)
           then one else zero)).
\end{lstlisting}

Both side conditions are facts about the Schnorr group, proved once in
\rocqfile[Utility/SchnorrOrder.v]{src/Utility/SchnorrOrder.v}: every
element other than $1$ of $\{v : v^q = 1 \bmod p\}$ has order $q$ for
prime $q$ (\rocq{prime_order}, by Bezout), so exponentiation with it is
injective (\rocq{pow_injective}), and $\mathbb{Z}_q$ has characteristic
$q$ (\rocq{natF_injective_below_q}). The instance then only checks $g
\ne 1$ and the bound, so \rocq{ballot_sound_ins},
\rocq{schulze_end_to_end_ins}, and \rocq{count_sound_ins} have no
hypothesis beyond the transcripts.

The zero-knowledge theorems \rocq{ballot_zk} and \rocq{tally_zk} have
the permutation form of Section~\ref{schn_proof}, with the simulators of
Section~\ref{sec:security}. Their proofs lift the library's per-protocol
simulators through the product distributions of the $m + 2m^2$ (resp.\
$O(m^3)$) proofs. The ballot simulator computes the $\delta$ ciphertexts
from the position ciphertexts as the verifier does.

Tallying is an inductive judgement, in the style of the evidence-carrying
counters of~\cite{10.1007/978-3-319-66107-0_26} and of the
\texttt{ECount} judgement of Haines et al.~\cite{haines2019verifiable}: \rocq{count}~\texttt{:~state -> Type} with rules
\rocq{ax}, \rocq{cvalid}, \rocq{cinvalid}, and \rocq{finish},
each carrying the invariant that the ballots seen are a permutation of
the accepted and the rejected ones. \rocq{compute_final_count} builds
a derivation for any list of ballots, so the extracted verifier returns
its verdict with the evidence. \rocq{count_sound} says that a
derivation ending in a true verdict, with rewinding transcripts, names a
Schulze winner of the rankings cast. We do not verify the honest
trustees' closure algorithm. \rocq{cert_ok_b} checks its plaintext
certificate before publication, and \rocq{verify_tally_sound} does not
depend on how a certificate was made.

The voter's client (\rocq{nizk_encrypt_ballot}), the verifier
(\rocq{nizk_verify_ballot}, \rocq{compute_final_count}), and the
trustees' prover (\rocq{nizk_generate_tally}) are extracted to OCaml
with the Fiat--Shamir hash (SHAKE-256) and the randomness as parameters
(Section~\ref{schn_proof}). The example program
(\rocqfile[Executable/Schulzecode]{src/Executable/Schulzecode}) prints the bulletin board of the
running example and walks the \rocq{count} derivation ballot by ballot.
The benchmark program (\rocqfile[Executable/SchulzeBenchcode]{src/Executable/SchulzeBenchcode}) runs the same
functions at the 2048-bit Helios parameters for any $m$ and $N$
(Section~\ref{sec:eval}). The Gentoo driver (\rocqfile[Executable/SchulzeGentoocode]{src/Executable/SchulzeGentoocode})
runs all of them on the published elections
(Section~\ref{sec:eval:gentoo}).
